\documentclass[12pt, a4paper]{article}

% --- PAQUETES ---
\usepackage[utf8]{inputenc} % Codificación de entrada
\usepackage[spanish]{babel} % Idioma español, para correcta separación de sílabas y títulos
\usepackage{geometry} % Para configurar los márgenes
\usepackage{amsmath} % Para entornos matemáticos avanzados
\usepackage{hyperref} % Para enlaces (si se necesitaran en el futuro)
\usepackage{titlesec} % Para personalizar los títulos de las secciones
\usepackage{enumitem} % Para personalizar las listas

% --- CONFIGURACIÓN DE PÁGINA Y TÍTULOS ---
\geometry{a4paper, margin=1in} % Márgenes de 2.54 cm en todos los lados

\titleformat{\section}
  {\normalfont\Large\bfseries\scshape\centering} % Formato: Normal, Grande, Negrita, Versalitas, Centrado
  {}
  {0em}
  {}[\titlerule] % Línea horizontal debajo del título de la sección
\titlespacing*{\section}{0pt}{2\baselineskip}{1.5\baselineskip}

\titleformat{\subsection}
  {\normalfont\large\bfseries} % Formato: Normal, Grande, Negrita
  {\thesubsection}
  {1em}
  {}

% --- DOCUMENTO ---
\begin{document}

\title{\bfseries Plan de Estudio de Física (10 Semanas)\\ \large Basado en "Física, 6ta Edición" de Wilson, Buffa y Lou}
\author{}
\date{}
\maketitle
\thispagestyle{empty}
\vspace{-2cm}

% --- SEMANA 1 ---
\section*{Semana 1: Fundamentos de la Mecánica y Cinemática}
\textit{El objetivo de esta semana es construir una base sólida en las unidades, la resolución de problemas y la descripción del movimiento en una dimensión.}
\subsection*{Capítulo 1: Medición y resolución de problemas}
\begin{itemize}[leftmargin=*]
    \item Unidades SI, conversión de unidades.
    \item Cifras significativas y análisis dimensional.
    \item Estrategias para la resolución de problemas.
\end{itemize}
\subsection*{Capítulo 2: Cinemática: descripción del movimiento}
\begin{itemize}[leftmargin=*]
    \item Distancia, rapidez, desplazamiento y velocidad.
    \item Aceleración y caída libre.
\end{itemize}

% --- SEMANA 2 ---
\section*{Semana 2: Movimiento Avanzado y Leyes de Newton}
\textit{Esta semana ampliarás tus conocimientos de cinemática a dos dimensiones y te introducirás en las causas del movimiento: las fuerzas.}
\subsection*{Capítulo 3: Movimiento en dos dimensiones}
\begin{itemize}[leftmargin=*]
    \item Componentes de vectores.
    \item Movimiento de proyectiles.
    \item Velocidad relativa.
\end{itemize}
\subsection*{Capítulo 4: Fuerza y movimiento}
\begin{itemize}[leftmargin=*]
    \item Las tres leyes de Newton.
    \item Diagramas de cuerpo libre y fricción.
\end{itemize}

% --- SEMANA 3 ---
\section*{Semana 3: Energía y Momento Lineal}
\textit{Aquí estudiarás dos de los conceptos más importantes y de mayor conservación en la física: la energía y el momento.}
\subsection*{Capítulo 5: Trabajo y energía}
\begin{itemize}[leftmargin=*]
    \item Trabajo, energía cinética y potencial.
    \item Teorema trabajo-energía y conservación de la energía.
    \item Potencia.
\end{itemize}
\subsection*{Capítulo 6: Cantidad de movimiento lineal y choques}
\begin{itemize}[leftmargin=*]
    \item Impulso y momento lineal.
    \item Conservación del momento.
    \item Choques elásticos e inelásticos.
\end{itemize}

% --- SEMANA 4 ---
\section*{Semana 4: Movimiento Rotacional y Gravitación}
\textit{Esta semana se centra en el movimiento de los cuerpos que giran y en la fuerza fundamental que gobierna los planetas y las estrellas.}
\subsection*{Capítulo 7: Movimiento circular y gravitacional}
\begin{itemize}[leftmargin=*]
    \item Movimiento circular uniforme y aceleración centrípeta.
    \item Ley de la gravitación de Newton.
\end{itemize}
\subsection*{Capítulo 8: Movimiento rotacional y equilibrio}
\begin{itemize}[leftmargin=*]
    \item Momento de fuerza (torque) y equilibrio.
    \item Momento de inercia y dinámica rotacional.
    \item Conservación del momento angular.
\end{itemize}

\newpage

% --- SEMANA 5 ---
\section*{Semana 5: Fluidos y Transición a la Termodinámica}
\textit{Concluirás el estudio de la mecánica con los fluidos y comenzarás a explorar el mundo de la temperatura y el calor.}
\subsection*{Capítulo 9: Sólidos y fluidos}
\begin{itemize}[leftmargin=*]
    \item Módulos de elasticidad.
    \item Presión, principio de Pascal y principio de Arquímedes.
    \item Dinámica de fluidos y ecuación de Bernoulli.
\end{itemize}
\subsection*{Capítulo 10: Temperatura y teoría cinética}
\begin{itemize}[leftmargin=*]
    \item Escalas de temperatura.
    \item Leyes de los gases ideales y expansión térmica.
\end{itemize}

% --- SEMANA 6 ---
\section*{Semana 6: Calor y Termodinámica}
\textit{Profundizarás en los conceptos de transferencia de energía en forma de calor y las leyes que gobiernan estos procesos.}
\subsection*{Capítulo 11: Calor}
\begin{itemize}[leftmargin=*]
    \item Calor específico y calorimetría.
    \item Cambios de fase y calor latente.
    \item Transferencia de calor (conducción, convección, radiación).
\end{itemize}
\subsection*{Capítulo 12: Termodinámica}
\begin{itemize}[leftmargin=*]
    \item Primera y segunda ley de la termodinámica.
    \item Máquinas de calor y entropía.
\end{itemize}

% --- SEMANA 7 ---
\section*{Semana 7: Vibraciones, Ondas y Sonido}
\textit{Esta semana te introducirás en el movimiento periódico y su propagación a través de los medios.}
\subsection*{Capítulo 13: Vibraciones y ondas}
\begin{itemize}[leftmargin=*]
    \item Movimiento armónico simple (MAS).
    \item Propiedades de las ondas y ondas estacionarias.
\end{itemize}
\subsection*{Capítulo 14: Sonido}
\begin{itemize}[leftmargin=*]
    \item Rapidez e intensidad del sonido.
    \item Efecto Doppler y fenómenos acústicos.
\end{itemize}

\newpage

% --- SEMANA 8 ---
\section*{Semana 8: Electrostática y Capacitancia}
\textit{Comienza la exploración de la electricidad y el magnetismo, empezando por las cargas en reposo.}
\subsection*{Capítulo 15: Cargas, fuerzas y campos eléctricos}
\begin{itemize}[leftmargin=*]
    \item Carga eléctrica y ley de Coulomb.
    \item Campo eléctrico y líneas de campo.
\end{itemize}
\subsection*{Capítulo 16: Potencial eléctrico, energía y capacitancia}
\begin{itemize}[leftmargin=*]
    \item Energía y potencial eléctrico.
    \item Capacitancia y dieléctricos.
\end{itemize}

% --- SEMANA 9 ---
\section*{Semana 9: Circuitos Eléctricos y Magnetismo}
\textit{Esta semana te enfocarás en el movimiento de las cargas (corriente) y los circuitos, y comenzarás con el magnetismo.}
\subsection*{Capítulo 17: Corriente eléctrica y resistencia}
\begin{itemize}[leftmargin=*]
    \item Baterías, corriente y ley de Ohm.
    \item Resistividad y potencia eléctrica.
\end{itemize}
\subsection*{Capítulo 18: Circuitos eléctricos básicos}
\begin{itemize}[leftmargin=*]
    \item Resistencias en serie y paralelo.
    \item Reglas de Kirchhoff.
\end{itemize}
\subsection*{Capítulo 19: Magnetismo}
\begin{itemize}[leftmargin=*]
    \item Fuerza magnética sobre cargas y corrientes.
    \item Campos magnéticos producidos por corrientes.
\end{itemize}

% --- SEMANA 10 ---
\section*{Semana 10: Inducción Electromagnética y Óptica}
\textit{La última semana unifica la electricidad y el magnetismo, para finalmente explorar la naturaleza y el comportamiento de la luz.}
\subsection*{Capítulo 20: Inducción y ondas electromagnéticas}
\begin{itemize}[leftmargin=*]
    \item Ley de Faraday de la inducción y ley de Lenz.
    \item Generadores, transformadores y ondas electromagnéticas.
\end{itemize}
\subsection*{Capítulo 22: Reflexión y refracción de la luz}
\begin{itemize}[leftmargin=*]
    \item Reflexión, refracción (ley de Snell).
    \item Dispersión y reflexión interna total.
\end{itemize}
\subsection*{Capítulo 23: Espejos y lentes}
\begin{itemize}[leftmargin=*]
    \item Espejos planos y esféricos.
    \item Lentes delgadas y ecuación del fabricante de lentes.
\end{itemize}

\newpage

% --- CONSEJOS ---
\section*{Consejos para el Éxito}
\begin{enumerate}[leftmargin=*]
    \item \textbf{Lectura Activa:} Antes de cada sesión de estudio, lee los capítulos correspondientes. No solo leas, intenta comprender los ejemplos resueltos.
    \item \textbf{Resolución de Problemas:} La física se aprende haciendo. Dedica la mayor parte de tu tiempo a resolver los ejercicios al final de cada capítulo. El libro tiene las respuestas a los ejercicios impares para que puedas verificar tu progreso.
    \item \textbf{Repaso Semanal:} Dedica un día del fin de semana para repasar los temas vistos y resolver algunos problemas adicionales para consolidar tu conocimiento.
    \item \textbf{No te quedes con dudas:} Si un concepto o problema te resulta difícil, busca ayuda. Puedes consultar el apéndice del libro, buscar recursos en línea o preguntar a un profesor o compañero.
\end{enumerate}

\end{document}